\documentclass[../DoAn.tex]{subfiles}
\begin{document}

Chương này tập trung trình bày chi tiết các giải pháp kỹ thuật nổi bật: sinh đề nhiều phiên bản kèm tọa độ OMR, nhận dạng OMR trên di động với phát hiện lỗi, và quản lý phiên làm việc và phân quyền.

\section{Giải pháp tự động sinh đề thi nhiều phiên bản và phiếu OMR kèm tọa độ nhận dạng}
\label{section:5.1}

\subsection{Dẫn dắt bài toán}
Trong Chương 2, mục \ref{subsection:2.2.6}, đồ án đã xác định yêu cầu tạo bài kiểm tra với khả năng tự động trộn câu hỏi thành nhiều phiên bản để hạn chế gian lận và sinh phiếu trả lời OMR tương ứng. Đây là bài toán khó vì mỗi phiên bản đề phải có đáp án đúng khác nhau do thứ tự câu hỏi và đáp án bị hoán vị, đồng thời tọa độ các ô trên phiếu phải được biết chính xác để ứng dụng di động có thể nhận dạng. Mục 4.3 chỉ mô tả kết quả đạt được; chi tiết giải pháp được trình bày tại đây.

\subsection{Giải pháp chi tiết}
Pipeline đồng bộ gồm năm giai đoạn chính, mỗi giai đoạn được thiết kế để chạy độc lập và có thể tái chạy khi cần sinh lại phiên bản đề.

\subsubsection{Giai đoạn 1 -- Sinh phiên bản đề và hoán vị đáp án}
Để hạn chế gian lận, mỗi học sinh trong cùng một lớp thi sẽ nhận được một phiên bản đề khác nhau. Thuật toán đảm bảo ba tính chất: (a) các phiên bản là hoán vị của nhau, (b) đáp án đúng của mỗi câu bị hoán vị theo thứ tự câu hỏi mới, (c) việc sinh là xác định nên có thể tái tạo kết quả.

Bước 1.1, lấy danh sách câu hỏi của đề thi gốc từ database. Mỗi câu hỏi bao gồm: nội dung, bốn lựa chọn (A, B, C, D), đáp án đúng, loại câu hỏi, độ khó.

Bước 1.2, áp dụng thuật toán Fisher--Yates shuffle để trộn thứ tự câu hỏi. Thuật toán duyệt mảng ngẫu nhiên và hoán đổi phần tử hiện tại với một phần tử ngẫu nhiên phía sau nó:
\begin{equation}
    \text{for } i = n-1 \text{ downto } 1: \quad j \gets \text{random}(0, i); \quad \text{swap}(A[i], A[j])
\end{equation}
Hạt giống (seed) của bộ sinh ngẫu nhiên được lưu lại để đảm bảo tính tái tạo.

Bước 1.3, với mỗi câu hỏi sau khi trộn, hoán vị thứ tự bốn lựa chọn (A, B, C, D) theo một hoán vị ngẫu nhiên khác nhau. Đáp án đúng được cập nhật tương ứng để phản ánh vị trí mới.

Bước 1.4, kết quả được lưu thành đối tượng \texttt{ExamVersion} gồm: \texttt{examId}, \texttt{versionIndex} (0 đến $N-1$), \texttt{seed} (hạt giống), \texttt{questionOrder} (danh sách ID câu theo thứ tự mới), \texttt{answerKey} (đáp án đúng dạng \{questionId: correctChoice\}), \texttt{shuffleMap} (ánh xạ vị trí câu gốc sang vị trí mới).

Pseudocode~\ref{alg:generateVersions} tóm tắt quy trình sinh phiên bản đề.

\begin{algorithm}[H]
\caption{Sinh phiên bản đề thi}
\label{alg:generateVersions}
\SetAlgoLined
\KwIn{Đề thi gốc $E$, số phiên bản $N$}
\KwOut{Danh sách $\mathcal{V} = \{V_0, V_1, \dots, V_{N-1}\}$}
 $\mathcal{V} \gets \emptyset$ \;
 \For{$i \gets 0$ \KwTo $N-1$}{
  $seed_i \gets \text{generateSeed}(E.id, i)$ \;
  $\text{rng} \gets \text{createRNG}(seed_i)$ \;
  $Q_{shuffled} \gets \text{fisherYatesShuffle}(E.questions, \text{rng})$ \;
  $A_{shuffled} \gets \emptyset$ \;
  \ForEach{$q \in Q_{shuffled}$}{
   $\pi_q \gets \text{randomPermutation}(4, \text{rng})$ \;
   $q.choices \gets \text{applyPermutation}(q.choices, \pi_q)$ \;
   $q.correctChoice \gets \pi_q^{-1}(q.correctChoice)$ \;
   $A_{shuffled}[q.id] \gets q.correctChoice$ \;
  }
  $V_i \gets \text{ExamVersion}(E.id, i, seed_i, Q_{shuffled}, A_{shuffled})$ \;
  $\mathcal{V} \gets \mathcal{V} \cup \{V_i\}$ \;
 }
 \Return{$\mathcal{V}$} \;
\end{algorithm}

\subsubsection{Giai đoạn 2 -- Sinh mã nguồn LaTeX cho đề và phiếu OMR}
Mỗi phiên bản đề cần hai file PDF: (a) file đề thi cho học sinh làm bài, (b) file phiếu trả lời OMR để chấm thi tự động. Hệ thống sử dụng thư viện AMC (Auto Multiple Choice) làm backend sinh LaTeX nhờ khả năng tạo cả hai loại file từ cùng một cấu hình.

Bước 2.1, module \texttt{amcSourceGenerator} nhận đầu vào là \texttt{ExamVersion} và cấu hình template (số câu, số lựa chọn, tiêu đề, logo trường). Generator duyệt danh sách câu hỏi theo thứ tự đã trộn và sinh mã LaTeX cho đề thi:
\begin{itemize}
  \item Môi trường \texttt{question} cho mỗi câu hỏi, bọc nội dung và bốn lựa chọn trong \texttt{choices}
  \item Macro \texttt{\textbackslash explication} cho giải thích (nếu có)
  \item Lệnh \texttt{\textbackslash multiplechoices} để đánh dấu đáp án đúng (không hiển thị cho học sinh)
\end{itemize}

Bước 2.2, module \texttt{amcAnswerSheetTexGenerator} sinh mã LaTeX cho phiếu trả lời OMR bằng môi trường \texttt{answerbox} của AMC. Cấu trúc phiếu gồm:
\begin{itemize}
  \item Vùng mã số sinh viên: $n$ cột, mỗi cột 4 bubble tương ứng ký tự số
  \item Vùng mã đề: $m$ cột, mỗi cột 4 bubble tương ứng ký tự chữ/số
  \item Vùng câu trả lời: ma trận $Q \times 4$ bubble, mỗi hàng một câu, mỗi cột một lựa chọn
  \item Khung viền bao quanh toàn bộ phiếu để phát hiện contour
\end{itemize}

Bước 2.3, mã LaTeX được ghi ra file \texttt{exam\_\textless version\_index\textgreater.tex} và \texttt{sheet\_\textless version\_index\textgreater.tex}. Mỗi file kèm theo header chứa metadata: mã đề, phiên bản, số câu, ngày tạo.

Pseudocode~\ref{alg:generateLatex} mô tả quy trình sinh mã nguồn.

\begin{algorithm}[H]
\caption{Sinh mã nguồn LaTeX}
\label{alg:generateLatex}
\SetAlgoLined
\KwIn{Danh sách phiên bản $\mathcal{V}$, cấu hình template $T$}
\KwOut{Danh sách file LaTeX}
 $\text{files} \gets \emptyset$ \;
 \ForEach{$V \in \mathcal{V}$}{
  $\text{examTex} \gets \text{amcExamHeader}(T)$ \;
  \ForEach{$q \in V.questionOrder$}{
   $\text{examTex} \gets \text{examTex} + \text{formatQuestion}(q, V.answerKey[q.id])$ \;
  }
  $\text{examTex} \gets \text{examTex} + \text{amcExamFooter}()$ \;
  $\text{writeFile}(\text{exam\_} + V.versionIndex + \text{.tex}, \text{examTex})$ \;

  $\text{sheetTex} \gets \text{amcSheetHeader}(T)$ \;
  $\text{sheetTex} \gets \text{sheetTex} + \text{formatStudentIdBox}(T)$ \;
  $\text{sheetTex} \gets \text{sheetTex} + \text{formatExamCodeBox}(T)$ \;
  $\text{sheetTex} \gets \text{sheetTex} + \text{formatAnswerGrid}(T, V.versionIndex)$ \;
  $\text{sheetTex} \gets \text{sheetTex} + \text{amcSheetFooter}()$ \;
  $\text{writeFile}(\text{sheet\_} + V.versionIndex + \text{.tex}, \text{sheetTex})$ \;
  $\text{files} \gets \text{files} \cup \{(\text{exam\_} + V.versionIndex, \text{sheet\_} + V.versionIndex)\}$ \;
 }
 \Return{$\text{files}$} \;
\end{algorithm}

\subsubsection{Giai đoạn 3 -- Trích xuất tọa độ các ô từ AMC}
Tọa độ chính xác của các bubble là thành phần then chốt để ứng dụng di động có thể nhận dạng. Hệ thống tận dụng chế độ \texttt{prepare} của AMC để sinh file tọa độ tự động.

Bước 3.1, biên dịch file LaTeX phiếu OMR bằng XeLaTeX trong môi trường WSL2 để đảm bảo tương thích với các gói LaTeX của AMC:
\begin{equation}
    \text{sheet\_}\textless i\textgreater.\text{tex} \xrightarrow{\text{XeLaTeX}} \text{sheet\_}\textless i\textgreater.\text{pdf}
\end{equation}

Bước 3.2, gọi lệnh AMC với chế độ \texttt{prepare} để phân tích cấu trúc phiếu:
\begin{equation}
    \texttt{auto-multiple-choice prepare --mode s --out-calage sheet\_\textless i\textgreater.tex}
\end{equation}
Chế độ \texttt{s} (separate) sinh file \texttt{sheet\_\textless i\textgreater.calage.xy} chứa tọa độ pixel của tất cả các vùng được đánh dấu trong phiếu.

Bước 3.3, module \texttt{amcCalageParser} đọc file \texttt{.calage.xy} theo định dạng AMC. File có cấu trúc:
\begin{itemize}
  \item Mỗi dòng tương ứng một vùng (box), có định dạng: \texttt{Box name x y width height}
  \item Các tọa độ được lưu ở đơn vị điểm (point) với độ phân giải nguồn
\end{itemize}

Bước 3.4, chuẩn hóa tọa độ về pixel ở 300 DPI:
\begin{equation}
    \text{pixel} = \text{point} \times \frac{300}{72} = \text{point} \times 4{,}1667
\end{equation}
Sau làm tròn đến số nguyên gần nhất, sai số dưới 0,5 pixel.

Bước 3.5, phân loại các box thành các nhóm chức năng: \texttt{studentIdBoxes}, \texttt{examCodeBoxes}, \texttt{answerBoxes}. Mỗi nhóm được sắp xếp theo thứ tự hàng và cột để tạo ánh xạ bubble -- câu hỏi -- lựa chọn.

Kết quả giai đoạn này là danh sách tọa độ bubble chuẩn hóa, được lưu vào cấu trúc \texttt{OMRTemplate} để ứng dụng Flutter sử dụng khi chấm thi.

\subsubsection{Giai đoạn 4 -- Tạo template OMR và ánh xạ câu hỏi}
Template OMR là cấu trúc trung tâm kết nối giữa backend (sinh đề) và frontend (nhận dạng OMR). Nó chứa tất cả thông tin không đổi của một dạng phiếu: vị trí các vùng đọc mã, cấu trúc lưới câu trả lời, và cách ánh xạ bubble sang câu hỏi.

Bước 4.1, từ kết quả giai đoạn 3, xây dựng \texttt{OMRTemplate} gồm các trường:
\begin{itemize}
  \item \texttt{examId}: liên kết với đề thi gốc
  \item \texttt{versionIndex}: phiên bản đề (0 đến $N-1$)
  \item \texttt{dpi}: độ phân giải chuẩn (300)
  \item \texttt{imageWidth}, \texttt{imageHeight}: kích thước ảnh chuẩn hóa
  \item \texttt{studentIdRegion}: tọa độ khung mã sinh viên, gồm danh sách \texttt{columns} (mỗi cột là list bubble với \{x, y, width, height\})
  \item \texttt{examCodeRegion}: tọa độ khung mã đề, cấu trúc tương tự
  \item \texttt{answerGrid}: danh sách hàng, mỗi hàng chứa \texttt{questionId}, \texttt{questionIndex}, danh sách \texttt{bubbles} (4 bubble tương ứng A, B, C, D)
\end{itemize}

Bước 4.2, xây dựng ánh xạ câu hỏi -- phiên bản -- đáp án. Vì thứ tự câu hỏi đã bị trộn, cần lưu ánh xạ để server biết đáp án đúng của câu thứ $k$ trong phiếu trả lời là gì:
\begin{equation}
    \text{answerKey}[k] = \text{correctChoice of question at position } k \text{ in version } V_i
\end{equation}

Bước 4.3, template được lưu vào database dưới dạng JSON trong collection \texttt{OMRTemplates}, kèm chỉ mục đơn (index) trên trường \texttt{\{examId, versionIndex\}} để truy vấn nhanh.

Bước 4.4, các file PDF (đề thi và phiếu OMR) được tải lên Cloudinary, lưu URL công khai. Backend lưu URL này trong \texttt{ExamVersion.pdfUrl} và \texttt{ExamVersion.sheetUrl} để giáo viên có thể in hoặc phát hành cho học sinh.

Pseudocode~\ref{alg:buildTemplate} mô tả quy trình tạo template OMR.

\begin{algorithm}[H]
\caption{Xây dựng template OMR}
\label{alg:buildTemplate}
\SetAlgoLined
\KwIn{File \texttt{.calage.xy} $\mathcal{B}$, phiên bản đề $V$}
\KwOut{Đối tượng \texttt{OMRTemplate}}
 $\text{boxes} \gets \text{parseCalage}(\mathcal{B})$ \;
 $\text{stdBoxes} \gets \text{normalizeTo300DPI}(\text{boxes})$ \;
 $\text{template} \gets \text{OMRTemplate}(examId=V.examId, versionIndex=V.versionIndex)$ \;
 $\text{template.studentIdRegion} \gets \text{extractRegion}(\text{stdBoxes}, \text{type=studentId})$ \;
 $\text{template.examCodeRegion} \gets \text{extractRegion}(\text{stdBoxes}, \text{type=examCode})$ \;
 $\text{template.answerGrid} \gets \text{buildAnswerGrid}(V.questionOrder, \text{stdBoxes})$ \;
 $\text{template.answerKey} \gets V.answerKey$ \;
 \Return{$\text{template}$} \;
\end{algorithm}

\subsubsection{Giai đoạn 5 -- Lưu trữ và phân phối đề thi}
Sau khi tất cả file đã được sinh và tải lên, hệ thống thực hiện lưu trữ có cấu trúc để phục vụ các nghiệp vụ sau: phát hành đề, chấm thi, thống kê.

Bước 5.1, lưu thông tin phiên bản vào database:
\begin{itemize}
  \item \texttt{ExamVersion}: lưu seed, thứ tự câu, đáp án, URL PDF, URL sheet
  \item \texttt{OMRTemplate}: lưu templateJson đầy đủ tọa độ bubble
  \item \texttt{Exam}: cập nhật trạng thái \texttt{published} và số phiên bản đã sinh
\end{itemize}

Bước 5.2, tạo quan hệ giữa học sinh và phiên bản đề. Hệ thống phân bổ ngẫu nhiên phiên bản cho từng học sinh trong lớp thi, lưu vào \texttt{StudentExamVersion}:
\begin{equation}
    \text{assignedVersion}_s = V_{s \bmod N}
\end{equation}
với $s$ là số thứ tự học sinh trong danh sách lớp. Phân bổ theo modulo đảm bảo các phiên bản được sử dụng đều nhau.

Bước 5.3, khi giáo viên phát hành đề, hệ thống gửi thông báo đến học sinh. Học sinh có thể xem đề của mình trên web hoặc in ra giấy. Phiếu OMR được in kèm theo đề hoặc phát hành riêng tùy theo cấu hình của trường.

Bước 5.4, khi giáo viên bắt đầu chấm thi bằng điện thoại, ứng dụng gọi API \texttt{GET /api/exams/\textless id\textgreater/template} để lấy \texttt{OMRTemplate} của phiên bản tương ứng với học sinh đang chấm. Template này được truyền vào \texttt{OMREngine.initialize()} để thiết lập tọa độ bubble cho toàn bộ pipeline nhận dạng.

\subsection{Kết quả đạt được}
Pipeline sinh tối đa 50 phiên bản đề kèm phiếu OMR trong dưới 2 phút cho đề 40 câu. Tọa độ các ô chính xác dưới 1 pixel ở 300 DPI. Đã thử nghiệm với 5 bài kiểm tra, hơn 100 phiên bản đề và 300 phiếu OMR; tất cả nhận dạng thành công.

\section{Giải pháp nhận dạng OMR trên thiết bị di động với khả năng phát hiện lỗi}
\label{section:5.2}

\subsection{Dẫn dắt bài toán}
Yêu cầu chấm thi bằng ứng dụng di động (mục \ref{subsection:2.2.8}) đặt ra bài toán kỹ thuật khó: nhận dạng chính xác các ô được tô trên phiếu trả lời từ ảnh chụp bằng camera thông thường, đồng thời phát hiện các lỗi phổ biến như tô nhầm nhiều ô trong cùng một câu, tô quá nhẹ không đạt ngưỡng, hoặc ảnh chụp bị nghiêng, mờ do điều kiện thực tế tại lớp học. Mục 4.4 đã trình bày kết quả kiểm thử; phần này trình bày chi tiết thuật toán.

\subsection{Giải pháp chi tiết}
Module \texttt{OMREngine} trong ứng dụng Flutter thực hiện pipeline nhận diện OMR gồm sáu giai đoạn chính. Mỗi giai đoạn được thiết kế độc lập, cho phép tinh chỉnh tham số mà không ảnh hưởng đến các bước khác.

\subsubsection{Giai đoạn 1 -- Tiền xử lý ảnh}
Ảnh chụp từ camera được đưa vào chuỗi xử lý chuẩn hóa nhằm giảm nhiễu và tăng độ tương phản trước khi phân tích. Các phép biến đổi được áp dụng theo thứ tự sau.

Bước 1.1, chuyển đổi không gian màu từ RGB sang grayscale để giảm kích thước dữ liệu và loại bỏ thông tin màu không cần thiết cho nhận dạng bubble:
\begin{equation}
    I_{gray}(x,y) = 0{,}299 \cdot I_R(x,y) + 0{,}587 \cdot I_G(x,y) + 0{,}114 \cdot I_B(x,y)
\end{equation}

Bước 1.2, lọc nhiễu bằng bộ lọc Gaussian kích thước $5 \times 5$ với độ lệch chuẩn $\sigma = 1{,}0$:
\begin{equation}
    G(x,y) = \frac{1}{2\pi\sigma^2} e^{-\frac{x^2+y^2}{2\sigma^2}}
\end{equation}
Phép tích chập $I_{blur} = I_{gray} * G$ giúp làm mịn nhiễu Gaussian và nhiễu muối tiêu, giữ lại cạnh chính của hình ảnh.

Bước 1.3, nhị phân hóa bằng phương pháp Otsu. Thuật toán tự động chọn ngưỡng $\tau_{otsu}$ tối ưu bằng cách tối đa hóa phương sai giữa hai lớp pixel nền và nội dung:
\begin{equation}
    \tau_{otsu} = \arg\max_{\tau} \bigl( \omega_0(\tau)\sigma_0^2(\tau) + \omega_1(\tau)\sigma_1^2(\tau) \bigr)
\end{equation}
trong đó $\omega_0$, $\omega_1$ là tỉ lệ pixel của hai lớp và $\sigma_0^2$, $\sigma_1^2$ là phương sai của từng lớp. Kết quả nhị phân:
\begin{equation}
    I_{bin}(x,y) = 
    \begin{cases}
        1 & \text{nếu } I_{blur}(x,y) > \tau_{otsu} \\
        0 & \text{nếu không}
    \end{cases}
\end{equation}

Bước 1.4, phép toán morphological opening và closing để loại bỏ nhiễu nhỏ và lấp lỗ trống trong vùng bubble:
\begin{align}
    I_{open} &= I_{bin} \circ SE = (I_{bin} \ominus SE) \oplus SE \\
    I_{clean} &= I_{open} \bullet SE = (I_{open} \oplus SE) \ominus SE
\end{align}
với $SE$ là phần tử cấu trúc hình elip $3 \times 3$.

Thuật toán tổng thể giai đoạn 1 được tóm tắt trong Pseudocode~\ref{alg:preprocess}.

\begin{algorithm}[H]
\caption{Tiền xử lý ảnh OMR}
\label{alg:preprocess}
\SetAlgoLined
\KwIn{Ảnh RGB $I_{RGB}$ kích thước $H \times W \times 3$}
\KwOut{Ảnh nhị phân đã làm sạch $I_{clean}$}
 $I_{gray} \gets \text{convertToGrayscale}(I_{RGB})$ \;
 $I_{blur} \gets \text{GaussianBlur}(I_{gray}, 5 \times 5, \sigma=1.0)$ \;
 $\tau \gets \text{OtsuThreshold}(I_{blur})$ \;
 $I_{bin} \gets \text{applyThreshold}(I_{blur}, \tau)$ \;
 $I_{clean} \gets \text{morphologicalClosing}(\text{morphologicalOpening}(I_{bin}))$ \;
 \Return{$I_{clean}$} \;
\end{algorithm}

\subsubsection{Giai đoạn 2 -- Phát hiện và hiệu chỉnh góc nghiêng}
Ảnh chụp từ camera thường bị nghiêng do góc cầm điện thoại không vuông góc với mặt phẳng phiếu. Để đảm bảo tọa độ bubble chính xác, thuật toán tự động phát hiện góc xoay và hiệu chỉnh.

Bước 2.1, phát hiện cạnh bằng toán tử Canny với ngưỡng thấp $t_{low}=50$, ngưỡng cao $t_{high}=150$, kích thước kernel Sobel $3 \times 3$:
\begin{equation}
    I_{edge} = \text{Canny}(I_{blur}, t_{low}, t_{high})
\end{equation}

Bước 2.2, biến đổi Hough xác định các đường thẳng chính của phiếu bằng cách tìm cực đại trong không gian accumulator:
\begin{equation}
    r = x \cos\theta + y \sin\theta, \quad \theta \in [0^\circ, 180^\circ]
\end{equation}
Hệ thống chọn hai đường thẳng có số phiếu bầu cao nhất và không gần song song với trục hoành để tính góc xoay $\alpha$. Nếu $|\alpha| < 1^\circ$, bước này bị bỏ qua để tiết kiệm thời gian.

Bước 2.3, xoay ảnh ngược chiều $-\alpha$ bằng phép affine transformation:
\begin{equation}
    I_{rot}(x',y') = I_{rot}\bigl( x\cos\alpha + y\sin\alpha,\; -x\sin\alpha + y\cos\alpha \bigr)
\end{equation}
Sau xoay, ảnh được cắt biên thừa để đảm bảo kích thước khớp với template OMR.

Pseudocode~\ref{alg:deskew} mô tả quy trình hiệu chỉnh góc nghiêng.

\begin{algorithm}[H]
\caption{Phát hiện và hiệu chỉnh góc nghiêng}
\label{alg:deskew}
\SetAlgoLined
\KwIn{Ảnh xám đã làm mờ $I_{blur}$}
\KwOut{Ảnh đã xoay thẳng $I_{rot}$}
 $I_{edge} \gets \text{Canny}(I_{blur}, 50, 150)$ \;
 $(r_1,\theta_1), (r_2,\theta_2) \gets \text{HoughLinesP}(I_{edge}, \text{topK}=2)$ \;
 $\alpha \gets \text{computeRotationAngle}(\theta_1, \theta_2)$ \;
 \uIf{$|\alpha| < 1^\circ$}{
  \Return{$I_{blur}$} \;
 }
 $M \gets \text{rotationMatrix}(\alpha, \text{center}=(W/2, H/2))$ \;
 $I_{rot} \gets \text{warpAffine}(I_{blur}, M, (W,H))$ \;
 \Return{$I_{rot}$} \;
\end{algorithm}

\subsubsection{Giai đoạn 3 -- Phát hiện vùng phiếu}
Sau khi ảnh đã thẳng, hệ thống cần xác định chính xác biên của phiếu trả lời để cắt vùng làm việc, loại bỏ nền và các vật thể nhiễu xung quanh.

Bước 3.1, tìm tất cả đường viền (contours) bằng hàm \texttt{findContours} trên ảnh nhị phân $I_{clean}$, sử dụng chuỗi lặp liền kề đơn giản với phương pháp trích xuất \texttt{RETR\_EXTERNAL}.

Bước 3.2, lọc contours theo tiêu chí diện tích. Chỉ giữ lại contours có diện tích $S$ nằm trong khoảng:
\begin{equation}
    S_{min} \le S \le S_{max}
\end{equation}
với $S_{min} = 0{,}15 \cdot H \cdot W$ và $S_{max} = 0{,}95 \cdot H \cdot W$. Các contours nhỏ hơn $S_{min}$ được loại bỏ vì có thể là nhiễu điểm hoặc ký tự in nhỏ; contours lớn hơn $S_{max}$ bị loại vì thường là viền khung hình ảnh.

Bước 3.3, từ các contours hợp lệ, chọn contour lớn nhất làm biên phiếu chính:
\begin{equation}
    C_{sheet} = \arg\max_{C_i \in \mathcal{C}_{valid}} \text{Area}(C_i)
\end{equation}

Bước 3.4, xấp xỉ $C_{sheet}$ thành tứ giác bằng phương pháp \texttt{approxPolyDP} với sai số tối đa $\varepsilon = 0{,}02 \cdot \text{perimeter}(C_{sheet})$. Nếu số đỉnh khác 4, thuật toán thử tìm tứ giác ngoài cùng của contours cha bằng cấu trúc phân cấp \texttt{RETR\_TREE}.

Bước 3.5, thực hiện phối cảnh biến đổi (perspective transform) để đưa tứ giác về hình chữ nhật chuẩn kích thước $W_{std} \times H_{std}$ (khớp với kích thước template OMR ở 300 DPI):
\begin{equation}
    I_{warp} = \text{perspectiveTransform}(I_{gray}, C_{sheet}, P_{std})
\end{equation}

Pseudocode~\ref{alg:detectSheet} tóm tắt quy trình phát hiện vùng phiếu.

\begin{algorithm}[H]
\caption{Phát hiện vùng phiếu trả lời}
\label{alg:detectSheet}
\SetAlgoLined
\KwIn{Ảnh nhị phân $I_{clean}$, ảnh gốc xám $I_{gray}$}
\KwOut{Ảnh phiếu đã chuẩn hóa $I_{warp}$}
 $\mathcal{C} \gets \text{findContours}(I_{clean}, \text{RETR\_EXTERNAL})$ \;
 $\mathcal{C}_{valid} \gets \{ C \in \mathcal{C} \mid S_{min} \le \text{Area}(C) \le S_{max} \}$ \;
 $C_{sheet} \gets \arg\max_{C \in \mathcal{C}_{valid}} \text{Area}(C)$ \;
 $P_{sheet} \gets \text{approxPolyDP}(C_{sheet}, \varepsilon=0.02 \cdot \text{peri}(C_{sheet}))$ \;
 \If{$|P_{sheet}| \ne 4$}{
  $P_{sheet} \gets \text{findLargestQuad}(P_{sheet})$ \;
 }
 $P_{std} \gets \text{standardRect}(W_{std}, H_{std})$ \;
 $I_{warp} \gets \text{warpPerspective}(I_{gray}, P_{sheet}, P_{std})$ \;
 \Return{$I_{warp}$} \;
\end{algorithm}

\subsubsection{Giai đoạn 4 -- Phát hiện bubble và đọc mã}
Vùng phiếu đã chuẩn hóa chứa hai loại thông tin cần trích xuất: (a) vùng cố định chứa mã đề và mã số sinh viên, (b) vùng câu trả lời gồm các bubble được sắp xếp theo lưới.

Bước 4.1, định vị vùng mã đề và mã số sinh viên dựa trên tọa độ tham chiếu từ template OMR. Mỗi vùng mã là một hình chữ nhật nhỏ được chia thành $n$ cột (số câu/số ký tự), mỗi cột chứa 4 bubble tương ứng 4 lựa chọn.

Bước 4.2, với mỗi bubble có tọa độ chuẩn $(x_i, y_i)$, hệ thống trích xuất vùng ảnh nhỏ quanh tâm bubble:
\begin{equation}
    R_i = I_{warp}\bigl[y_i - r : y_i + r,\; x_i - r : x_i + r\bigr]
\end{equation}
với bán kính $r = 12$ pixel ở 300 DPI.

Bước 4.3, đánh giá mức độ tô của bubble bằng công thức tỉ lệ tô đậm (fill ratio):
\begin{equation}
    f_i = \frac{\text{Area}(R_i \cap I_{bin})}{\text{Area}(R_i)}
\end{equation}
trong đó $I_{bin}$ là ảnh nhị phân của vùng $R_i$. Giá trị $f_i \in [0,1]$ thể hiện tỉ lệ pixel đen bên trong vùng bubble.

Bước 4.4, xác định trạng thái bubble:
\begin{equation}
    \text{status}_i = 
    \begin{cases}
        \text{FILLED} & \text{nếu } f_i \ge \theta \\
        \text{EMPTY} & \text{nếu } f_i < \theta_{light} \\
        \text{LIGHT} & \text{nếu } \theta_{light} \le f_i < \theta
    \end{cases}
\end{equation}
với ngưỡng chính $\theta = 0{,}4$ và ngưỡng tô nhẹ $\theta_{light} = 0{,}15$.

Bước 4.5, đọc mã số sinh viên và mã đề bằng cách chọn bubble có $f_i$ lớn nhất trong mỗi cột, trả về ký tự tương ứng (A, B, C, D). Nếu không có bubble nào vượt $\theta$, cột đó được đánh dấu là không đọc được.

Pseudocode~\ref{alg:readBubbles} mô tả quy trình đọc bubble và mã.

\begin{algorithm}[H]
\caption{Phát hiện bubble và đọc mã}
\label{alg:readBubbles}
\SetAlgoLined
\KwIn{Ảnh đã chuẩn hóa $I_{warp}$, danh sách bubble $\mathcal{B} = \{(x_i,y_i)\}$}
\KwOut{Map mã đề, Map mã sinh viên, Map đáp án}
 $\text{codes} \gets \emptyset$; $\text{answers} \gets \emptyset$ \;
 \ForEach{$\text{codeRegion} \in \text{template.codeRegions}$}{
  \ForEach{$\text{col} \in \text{codeRegion.columns}$}{
   $b_{best} \gets \arg\max_{b \in \text{col}} f(b)$ \;
   \If{$f(b_{best}) \ge \theta$}{
    $\text{codes}[\text{region}] \gets \text{col.indexOf}(b_{best})$ \;
   }
  }
 }
 \ForEach{$\text{row} \in \text{template.answerRows}$}{
  $b_{best} \gets \arg\max_{b \in \text{row}} f(b)$ \;
  \If{$f(b_{best}) \ge \theta$}{
   $\text{answers}[\text{row.questionNo}] \gets \text{row.choices}[b_{best}]$ \;
  } \ElseIf{$\exists b \in \text{row} : f(b) \ge \theta_{light}$}{
   $\text{answers}[\text{row.questionNo}] \gets \text{FLAG\_LIGHT}$ \;
  }
  \If{$|\{ b \in \text{row} : f(b) \ge \theta \}| > 1$}{
   $\text{answers}[\text{row.questionNo}] \gets \text{FLAG\_MULTI}$ \;
  }
 }
 \Return{$\text{codes}, \text{answers}$} \;
\end{algorithm}

\subsubsection{Giai đoạn 5 -- Đánh giá câu trả lời và phát hiện lỗi}
Sau khi trích xuất đáp án thô từ phiếu, hệ thống thực hiện đánh giá so với đáp án chuẩn của phiên bản đề tương ứng và phát hiện các trường hợp bất thường.

Bước 5.1, so sánh đáp án đọc được với đáp án đúng từ \texttt{ExamVersion.answerKey}:
\begin{equation}
    \text{score} = \sum_{q=1}^{Q} \mathbb{1}\bigl( \text{readAns}_q = \text{correctAns}_q \bigr) \times \frac{10}{Q}
\end{equation}
với $Q$ là tổng số câu hỏi.

Bước 5.2, phát hiện lỗi tô nhầm. Với mỗi câu $q$, đếm số bubble thỏa $f_i \ge \theta$:
\begin{equation}
    n_{marked}(q) = \bigl| \{ b \in \text{row}_q \mid f(b) \ge \theta \} \bigr|
\end{equation}
Nếu $n_{marked}(q) > 1$, hệ thống ghi nhận lỗi \texttt{MULTI\_MARK} và không tính điểm câu đó, đồng thời lưu danh sách bubble được tô để hiển thị lại cho giáo viên.

Bước 5.3, phát hiện lỗi tô nhẹ. Nếu không có bubble nào vượt ngưỡng $\theta$ nhưng tồn tại bubble với $\theta_{light} \le f_i < \theta$, ghi nhận lỗi \texttt{LIGHT\_MARK}. Câu này được đánh dấu cần kiểm tra thủ công.

Bước 5.4, phát hiện lỗi bỏ trống. Nếu $n_{marked}(q) = 0$ và không có bubble nào rơi vào khoảng $[\theta_{light}, \theta)$, ghi nhận lỗi \texttt{BANK\_Graded}.

Bước 5.5, tổng hợp kết quả thành đối tượng \texttt{OMRResult} gồm: điểm số, danh sách đáp án đọc được, danh sách lỗi theo câu, ảnh đã xử lý có vẽ biên bubble để hiển thị minh họa. Giáo viên có thể chấp nhận kết quả hoặc điều chỉnh điểm thủ công trên giao diện web.

Pseudocode~\ref{alg:evaluate} mô tả quy trình đánh giá và phát hiện lỗi.

\begin{algorithm}[H]
\caption{Đánh giá câu trả lời và phát hiện lỗi}
\label{alg:evaluate}
\SetAlgoLined
\KwIn{Đáp án đọc được $\text{readAns}$, đáp án chuẩn $\text{correctAns}$, bảng $\mathcal{B}$}
\KwOut{Điểm số, danh sách lỗi}
 $\text{score} \gets 0$; $\text{errors} \gets \emptyset$ \;
 \ForEach{$q \in [1,Q]$}{
  $n_{marked} \gets |\{ b \in \text{row}_q : f(b) \ge \theta \}|$ \;
  $n_{light} \gets |\{ b \in \text{row}_q : \theta_{light} \le f(b) < \theta \}|$ \;
  \uIf{$n_{marked} = 1$}{
   \If{$\text{readAns}_q = \text{correctAns}_q$}{
    $\text{score} \gets \text{score} + \frac{10}{Q}$ \;
   }
  }
  \uIf{$n_{marked} > 1$}{
   $\text{errors}[q] \gets \text{MULTI\_MARK}$ \;
  }
  \ElseIf{$n_{marked} = 0 \land n_{light} > 0$}{
   $\text{errors}[q] \gets \text{LIGHT\_MARK}$ \;
  }
  \ElseIf{$n_{marked} = 0 \land n_{light} = 0$}{
   $\text{errors}[q] \gets \text{BANK\_Graded}$ \;
  }
 }
 \Return{$\text{score}, \text{errors}$} \;
\end{algorithm}

\subsubsection{Giai đoạn 6 -- Tính điểm và lưu kết quả}
Kết quả cuối cùng của pipeline là đối tượng \texttt{OMRResult} chứa: (a) điểm số tổng và điểm chi tiết từng câu, (b) mã đề và mã sinh viên đọc được, (c) danh sách lỗi kèm mức độ nghiêm trọng, (d) ảnh đã xử lý có annotation biên bubble và mã QR/nhận dạng. Ứng dụng hiển thị kết quả tạm thời cho giáo viên xem trước trong 5 giây, sau đó tự động gửi lên server qua API \texttt{POST /api/submissions/omr}. Server thực hiện kiểm tra quyền sở hữu đề thi, lưu vào MongoDB và kích hoạt thông báo cho phụ huynh nếu có bật tính năng.

\subsection{Kết quả đạt được}
Thuật toán đạt độ chính xác 96,2\% trên 300 phiếu thử nghiệm thực tế. Thời gian xử lý trung bình 2,3 giây/phiếu trên điện thoại Android tầm trung. Phát hiện 100\% trường hợp tô nhầm và tô nhẹ. Hệ thống đã xử lý thành công hơn 900 phiếu trong môi trường lớp học thực tế với đủ điều kiện ánh sáng và góc chụp khác nhau.
\section{Giải pháp quản lý phiên làm việc và phân quyền theo vai trò trên đa nền tảng}
\label{section:5.4}

\subsection{Dẫn dắt bài toán}
Hệ thống Xây dựng hệ thống quản lý ngân hàng câu hỏi và chấm thi trắc nghiệm tự động sử dụng công nghệ OMR phục vụ năm vai trò người dùng khác nhau (admin, school-admin, teacher, student, parent) trên cả ba nền tảng (web, di động, API). Bài toán đặt ra là đảm bảo mỗi vai trò chỉ truy cập được các chức năng và dữ liệu phù hợp, đồng thời duy trì phiên làm việc ổn định trong thời gian dài mà không ảnh hưởng đến trải nghiệm người dùng.

\subsection{Giải pháp chi tiết}
Xác thực dùng JWT với cặp access token (15 phút) và refresh token (7 ngày). Khi access token hết hạn, ứng dụng tự gọi API refresh-token. Refresh token lưu trong cơ sở dữ liệu với thông tin thiết bị để thu hồi khi cần.

Phân quyền hai lớp: lớp định tuyến (middleware auth kiểm tra JWT, middleware permission kiểm tra vai trò) và lớp nghiệp vụ (service kiểm tra quyền sở hữu dữ liệu). Trên web, route guard trong React Router chuyển hướng theo vai trò (admin → /admin, school-admin → /school, teacher → /, student → /my-scores), sidebar lọc menu theo vai trò. Trên di động, BLoC auth quản lý điều hướng qua \texttt{go\_router} với route có bảo vệ.

\subsection{Kết quả đạt được}
100\% endpoint API được bảo vệ bằng kiểm tra quyền. Thời gian phản hồi tăng không đáng kể (khoảng 5ms). Cơ chế refresh tự động giảm phiền toái cho người dùng.

Chương 5 đã trình bày ba giải pháp kỹ thuật nổi bật: pipeline sinh đề-phiếu OMR với tọa độ chính xác, thuật toán OMR trên di động đạt 96,2\%, và quản lý phiên làm việc và phân quyền. Chương 6 sẽ tổng kết, so sánh với sản phẩm tương tự và đề xuất hướng phát triển.

\end{document}